% 다른 원고로 복사하여 사용할 수 있는 Related Work 초안.
% 현재 main.tex에서는 불러오지 않는다.
\section{관련 연구}
\label{sec:related_work_draft}

본 연구와 직접 연결되는 선행연구는 크게 세 가지로 구분된다. 첫째는
스크린샷으로부터 웹페이지의 내용과 요소 위치를 학습하는 Web/GUI
사전학습이고, 둘째는 GUI agent를 위한 지각ㆍ그라운딩ㆍ행동 데이터의
결합이며, 셋째는 caption의 세부 정보와 포괄성이 시각--언어 학습에
미치는 영향이다. 도구 호출이나 HTML 기반 agent는 스크린샷 기반
접근법의 동기를 설명하는 배경이지만, 본 연구가 비교하는 변수는 아니므로
Introduction에서 간략히 다루고 여기서는 제외한다.

\subsection{스크린샷 기반 Web UI 이해}

웹페이지는 자연 이미지와 달리 한 화면 안에 조밀한 텍스트, 수치,
상태, 아이콘, 입력 요소와 공간적 구조를 함께 포함한다. WebSRC는
웹페이지 질의응답에서 텍스트 내용뿐 아니라 문서 구조를 함께 사용하는
문제를 제시했고 \citep{chen2021websrc}, VisualWebBench는 웹페이지
이해를 captioning, QA, OCR, element grounding과 action 관련 능력으로
분해하여 멀티모달 모델을 평가했다 \citep{liu2024visualwebbench}.
이러한 연구는 웹 이해가 페이지의 전역 의미만으로 환원되지 않으며,
세부 텍스트와 요소 수준의 위치 및 기능을 함께 요구함을 보여준다.

이 능력을 학습하기 위해 기존 연구는 서로 다른 과제의 supervision을
혼합한다. MultiUI는 webpage caption, QA, heading 및 element OCR,
grounding과 action prediction을 합성하고, 이를 이용한 GUI knowledge
learning이 웹페이지뿐 아니라 문서ㆍ차트 이해에도 전이됨을 보고했다
\citep{liu2025multiui}. SeeClick은 화면 요소와 좌표의 대응을 학습하는
GUI grounding continual pre-training을 수행한 뒤 downstream agent에
적응시켰다 \citep{cheng2024seeclick}. GUI-Course 역시 화면 텍스트,
레이아웃과 좌표를 포함한 대규모 supervision을 사용하여 OCR,
grounding 및 navigation 성능의 변화를 분석했다
\citep{chen2025guicourse}. 이 연구들은 웹 전용 perception
supervision의 유용성을 보이지만, caption, QA, OCR과 grounding의
데이터 규모 및 출력 형식이 함께 달라지므로 description에 보존되는
정보량의 효과를 독립적으로 설명하기는 어렵다.

\subsection{GUI Agent를 위한 Perception과 행동 학습}

최근의 native GUI agent는 지각, 그라운딩, 행동 예측과 trajectory를
단계적으로 연결하거나 하나의 학습 mixture로 결합한다. UI-TARS는
element description, dense caption과 두 화면 사이의 state transition
caption을 perception data로 사용하고, 이후 action 및 trajectory
supervision을 적용한다 \citep{qin2025uitars}. MolmoWeb은 대규모
screenshot QA와 referring-expression grounding을 trajectory 중심의
SFT mixture에 포함하여 화면 이해와 웹 행동을 함께 학습한다
\citep{gupta2026molmoweb}. Aguvis와 InfiGUIAgent도 grounding 중심의
기초 단계와 planning 또는 expectation--reflection 학습을 구분한다
\citep{xu2025aguvis,liu2026infiguiagent}.

이러한 recipe는 agent 행동 이전에 화면의 내용과 요소를 이해하는
supervision이 필요하다는 공통된 설계 원칙을 보여준다. 그러나 각 연구는
caption, QA, grounding, action 및 trajectory의 비율과 데이터 출처가
서로 다르며, 초기 description의 정보 범위를 독립 변수로 비교하지
않는다. 따라서 최종 agent 성능 차이가 description 자체에서 비롯된
것인지, 다른 perception task나 후속 trajectory에서 비롯된 것인지
분리하기 어렵다. 본 연구는 grounding과 후속 IT, action warm-up,
trajectory 데이터를 고정하고, 사전학습 description의 information
budget만 달리하여 이 간극을 다룬다.

\subsection{Caption의 세부 정보와 Information Coverage}

일반 이미지 기반 MLLM 연구에서도 짧은 caption이 이미지와 언어의
기본적인 정렬에는 유용하지만, 세부 속성이나 관계를 충분히 전달하지
못할 수 있다는 문제가 제기되었다. LLaVA는 대규모 image--caption
pair를 이용해 시각 표현과 언어 모델을 연결한 뒤 instruction tuning을
수행하는 학습 recipe를 확립했다 \citep{liu2023llava}. ShareGPT4V는
일반 이미지의 상세한 description을 합성하고 기존 caption을 더 풍부한
설명으로 대체했을 때 여러 MLLM에서 성능이 향상됨을 보고했다
\citep{chen2024sharegpt4v}. 이 결과는 시각--언어 정렬에서 단순히
문장이 긴지보다 이미지에 근거한 세부 내용을 얼마나 정확하게 제공하는지가
중요할 수 있음을 시사한다.

Caption의 품질은 정확성뿐 아니라 density와 coverage의 관점에서도
분석되어 왔다. DAC는 caption--image alignment와 caption density를
구분하고, 정렬된 caption이라도 이미지의 객체, 속성, 상태와 관계를
충분히 언급하지 않으면 제한된 supervision이 될 수 있음을 보였다
\citep{doveh2023dac}. LoTLIP은 긴 텍스트를 활용하는 image--text
pre-training에서 긴 description을 그대로 하나의 target으로 사용하는
것만으로 충분하지 않으며, 세부 내용과 전역 의미의 학습 방법도 중요함을
보였다 \citep{wu2024lotlip}. 또한 Long Story Short는 caption의 길이와
세부 정보량을 변화시킨 실험에서 더 많은 정보가 항상 단조로운 향상으로
이어지지는 않으며, 중복과 최적화 부담이 개입할 수 있음을 보고했다
\citep{salazar2026longstory}.

그러나 이러한 연구의 대부분은 자연 이미지를 대상으로 하며, 웹페이지에
특유한 가시 텍스트, 선택 상태, 요소 기능, 좌표와 affordance가 caption에
포함되는 정도를 통제하지 않는다. 반대로 Web/GUI 연구는 dense caption을
사용하더라도 다른 QA, grounding 및 agent data와 함께 학습하여 caption
coverage의 독립 효과를 제시하지 않는다. 본 연구는 동일한 웹 스크린샷과
동일한 구조화 evidence로부터 Info-S, Info-M, Info-L을 생성한다.
세 조건은 각각 salient, representative, near-complete information
budget을 제공하며, 동일한 grounding 및 후속 학습을 적용한다. 따라서
본 연구의 초점은 ``긴 caption이 항상 우수한가''가 아니라, text-rich
웹 화면에서 어떤 수준의 evidence verbalization이 웹 이해와 최종 agent
적응에 효과적인가를 검증하는 데 있다.
